\subsection{The replacement premium connects the old and new groups}
\label{sec:insurance-modification-values}

A qualifying modification changes both the removed obligation and the new
contract. Paragraph~77 prevents the insurer from choosing the new opening
balance simply by copying the old one. For a contract derecognised under
paragraph~72, determine the premium that would be charged for an equivalent
new contract at the modification date, then subtract any additional premium
actually charged for the modification. This difference adjusts the old
group's CSM together with the removed fulfilment cash flows, and is treated
as a premium received when measuring the new contract
\citep[IFRS~17.72--77]{eu20231803}.

Take a general-measurement group with FCF 100, comprising expected claims
90 and risk adjustment 10, and CSM 30. Cash is 200, so liability 130 leaves
equity 70. A legally effective modification substantially changes one
contract's boundary and requires a new contract under paragraph~72. Its old
allocated FCF is 40, comprising claims 36 and risk adjustment 4. Assume zero
discount rates and no intervening service, tax, FX or OCI. The independently
supplied equivalent new premium is 60, of which 10 is actually charged and
received as an additional premium. The new obligation has FCF 52, comprising
claims 48 and risk adjustment 4. These values are case inputs; they are not
inferred from the old carrying amount.

The deemed premium is \(60-10=50\). Removing FCF 40 while transferring
value 50 reduces the old group's CSM by 10. Its remaining FCF is 60 and
CSM 20, giving liability 80. The new contract's total consideration of 60
less FCF 52 gives CSM 8. One accounting sequence debits the old liability
50 and credits the new liability 50, then debits cash 10 and credits the
new liability 10 for the actual receipt. Cash 210 equals old liability 80,
new liability 60 and unchanged equity 70. Recording a second cash receipt
for the deemed 50 would invent liquidity.

If the old group previously had three equal remaining coverage units and
the removed contract contributed one, it now has two; the new group has its
own service assessment. A proportional removal of 40\% of the old CSM
would leave 18 instead of 20 and would disregard the equivalent-premium
rule. Over the remaining service, assume actual claims equal the supplied
estimates and all remaining risk adjustment and CSM are released. The old
group produces profit \(80-54=26\); the new group produces
\(60-48=12\). Final cash \(210-54-48=108\) equals opening equity 70
plus profit 38.

The same removed FCF with an actual third-party transfer premium of 50
instead reduces the old liability by 50 and pays cash 50; there is no new
issued contract in this separate case. Cash 150, liability 80 and equity
70 reconcile. If the original CSM were only 6, either adjustment requiring
10 would consume it and recognise service loss 4. For the modification,
old liability 60 plus new liability 60 would then leave equity 90 from
cash 210, compared with opening equity \(200-106=94\). The resulting
loss component 4 is tracked within the old liability, not added again.
Paragraphs~44(c), 48--50 and 77 govern that limit
\citep[IFRS~17.44,48--50,77]{eu20231803}.

Where the removed contract carries allocated accumulated OCI, paragraph~91
adds a separate presentation decision. Amounts arising under the systematic
allocation in 88(b) are reclassified to profit on the specified transfer or
derecognition; amounts arising under the held-underlying-item policy in
89(b) are not. Preserve that origin and the allocated amount before closing
the old contract. The no-OCI example above does not eliminate this step
\citep[IFRS~17.91]{eu20231803}.
